Digitale data is kwetsbaar. Je kan het per ongeluk overschrijven of weggooien. Het kan gelockt worden door een crypto-locker of gestolen worden door een kwaadwillende. Data kan ook verloren gaan door diefstal, brand of het defect gaan van opslagsystemen. Het terug krijgen van de juiste data kan een heleboel werk opleveren als we geen goede backups hebben van de data of zelfs onmogelijk blijken te zijn.

Moderne systemen zoals smartphones zijn vaak verbonden met cloud-opslagsystemen waarbij er een automatische backup gemaakt wordt van de belangrijkste data. Dat geldt voor Linux systemen vaak niet. Als we onze data veilig willen stellen dan zullen we dat zelf moeten doen en moeten we erover nadenken hoe we dat gaan doen. Dit document bevat een aantal commando's die je kunnen helpen bij het veiligstellen van jou data of de data van gebruikers waarvoor je verantwoordelijk bent.

